% !PRINTED 298
\weberpage{298}
%<<CONT
divisor $Q'$ of $Q$ which is different both from $Q$ itself and from the unit group. But this divisor $Q'$ must be a normal divisor of $P$, because, if $\pi$ is contained in $Q$ and $\gamma$ in $P$, then $\gamma^{-1}\pi\gamma$, which is of the same degree as $\pi$, must likewise be contained in $Q$, and hence belongs to $Q'$ if $\pi$ belongs to $Q'$. But since, as remarked above, $Q$ can have no proper divisor that is greater than the unit group and at the same time a normal divisor of $P$, we must have $\lambda = 1$.

\begin{center}
\S.\,76.\\
\gesperrt{Representation of the Abelian group $Q$.}
\end{center}

The considerations of the preceding section have shown that a metacyclic group $P$, which can occur as the \gesperrt{Galois} group of an irreducible primitive equation $f(x) = 0$ of degree $n = p^k$, must contain an \gesperrt{Abelian group} $Q$ as a normal divisor, whose degree is $p^k$, and which, apart from the unit element, contains only elements of degree $p$.

By \S.\,9 this group $Q$ can be represented by a basis consisting of $k$ elements $A_1, A_2, \ldots, A_k$ of degree $p$, so that every element of $Q$ takes the form:
\weq{1}{A_1^{z_1} A_2^{z_2} \ldots A_k^{z_k},}
and here $z_1, z_2, \ldots, z_k$ run, independently of one another, each through a complete system of residues to the modulus $p$. Now, in order to find the corresponding permutation group, we start from an arbitrary root $x$ of $f(x)$, denote the root into which $x$ passes under the permutation (1) by
\weq{2}{[z_1, z_2, \ldots, z_k],}
and stipulate that this symbol shall not change its meaning when the numbers $z_1, z_2, \ldots, z_k$ are altered arbitrarily by multiples of $p$. To the root $x$ from which we started there then belongs the symbol $[0, 0, \ldots, 0]$, and by the symbol (2) every root of $f(x)$ is represented once and only once.

If $x$ passes into $x'$ under a permutation $A'$, and $x'$ into $x''$ under $A''$, then $x$ passes into $x''$ under the permutation $A' A''$. From this it follows, however, that $[z_1, z_2, \ldots, z_k]$, under the permutation
\weq{3}{A = A_1^{\alpha_1} A_2^{\alpha_2} \ldots A_k^{\alpha_k}}
%CONT>>
